\section{Why the computation is cheap}\label{app:cost}

This appendix compares the computation of Section~\ref{sec:computation} with two earlier ones. One is the KSS
Matlab package, in its routine \texttt{leave\_out\_COMPLETE}. The other is \fevc's structured route,
\texttt{inference(highrank)}, which is behind the standard errors in the \fevc{} paper. Table~\ref{tab:cost} counts linear solves beyond those of
the point estimate. The counts for the KSS Matlab package come from reading its code (the simulation routine
is the same in the package's GitHub version of October 2024, commit \texttt{8b957ffe}; other options of the
package, such as smoothing on a subsample, are not counted), those for \fevc{} from its documentation of November
2026 (\texttt{standard\_errors.tex}, repository commit \texttt{53947b43}); we did not time either. The column for
this note is the standard error alone, without the diagnostics of Section~\ref{sec:evidence}. The comparison is
kept as it was made before the implementation. \fevc{} 0.6.0 removed the structured route and the development
default mentioned below, and \code{inference(highrank)} now selects the normal intervals of the method of this note
(Section~\ref{sec:impl-diff}, item~10). With its default eight draws the pass takes 35 solves, and about 90 with
the leading directions that the default computes for every component, plus 400 for the offset pass with controls
under the default \code{nuisance(joint)} (Section~\ref{sec:computation});
Section~\ref{kss:sec:cost} compares the released package with the two KSS codes.

\begin{table}[H]
\centering\small
\caption{Linear solves for the standard errors of the variance components, beyond the point estimate}\label{tab:cost}
\begin{tabular}{@{}p{0.27\textwidth}p{0.2\textwidth}p{0.2\textwidth}p{0.22\textwidth}@{}}
\toprule
 & KSS Matlab package & \fevc{}, structured route & This note \\
\midrule
Error variances & smoothing of $\hat\sigma_i^2$ over all observations (lowess), per component; no solves & 2{,}048 solves to fit a variance model & closed form for stayers, two small regressions for movers; no solves \\
\addlinespace
Influences $\hat m$ & 2 per component & 3 & 3 \\
\addlinespace
Double sum & 1{,}000 draws per component, one number per draw & 1{,}000 draws, one number per component and draw & 16 draws, four solves each: 64 \\
\addlinespace
Dominance diagnostic & eigenvalue routines, on request & 8{,}850 & a bound from the same draws; no solves \\
\midrule
Total & about 3{,}000 for three components & 11{,}901 & 67 \\
\bottomrule
\multicolumn{4}{@{}p{0.93\textwidth}@{}}{\rule{0pt}{2.6ex}KSS Matlab package: each of its draws also solves a second,
block-diagonal system. \fevc's development default of the time (the leave-out residual standard error, never
released) has the same pattern of solves as this note (three and four per draw; 35 with its 8 draws) with different
variance estimates.}
\end{tabular}
\end{table}

Four things account for the difference.

\emph{One vector per draw instead of one number.} The variance of $\hat\theta$ contains the double sum
$\sum\sum C_{i\ell}^2\sigma_i^2\sigma_\ell^2$. The two earlier computations draw an error vector $\zeta$ with the
fitted variances, compute the single number $\zeta'C\zeta$ (one solve), and take the sample variance of that
number across draws. With $r_{\rm eff}=(\sum_j\lambda_j^2)^2/\sum_j\lambda_j^4$ the effective number of
directions of the matrix with entries $\st_iC_{i\ell}\st_\ell$ and $\lambda_j$ its eigenvalues, that sample
variance has relative variance $2/(R-1)+12/(R\,r_{\rm eff})$ from $R$ draws: at least $2/(R-1)$ whatever the
design, hence the thousand draws. Section~\ref{sec:computation} also computes the vector $C\zeta$, at one more
solve per component: $\sum_i\st_i^2(C\zeta)_i^2$ is an average over millions of rows, and one draw already has
relative variance $2/r_{\rm eff}$. On a large,
well-connected panel $r_{\rm eff}$ is in the thousands, so a handful of draws is enough. On the Veneto panel it
is between about \R{real.reff.min} and \R{real.reff.max} for the four components (estimated from 32 draws), and
runs with different sets of draws, and a run with 64 draws, differ by at most \R{real.run_maxdiff} in the
standard error.

\emph{The same draws serve every component.} One draw costs one solve for its residual and one for each of
three components, and gives the double sums of all four components and of their covariances. The KSS Matlab package simulates each component separately.

\emph{The error variances need no solves.} A stayer's variance estimate is a sample variance, and the model for
movers is two regressions with eleven coefficients on quantities the point estimate has already produced.
The KSS Matlab package smooths the leave-out estimates over all observations, once per component; the
structured route of \fevc{} fits its variance model from 2{,}048 further solves.

\emph{A dominance diagnostic comes with the draws.} With $u_d$ the vector with entries $\st_i(C\zeta^{(d)})_i$,
the average of $(u_d'u_e)^2$ over pairs of different draws is unbiased for $\sum_j\lambda_j^4$, given the data,
and $(\sum_j\lambda_j^4)^{1/2}/\hat\V[\hat\theta]$ is an upper bound for the ratio
$\lambda_{\max}^2/\hat\V[\hat\theta]$. This is an estimated bound, and it is for the matrix with the fitted
variances, not for the one with the true variances in Proposition~\ref{prop:conc}. On the Veneto panel the estimated bound from the 16 draws of the standard error is between
\R{real8.d2.min} and \R{real8.d2.max} for the four components, in two sets of draws. The ratio itself, by power
iteration, is between \R{real.lambda.min} and \R{real.lambda.max} (Section~\ref{sec:evidence}; 20 iterations per
component, each applying the weights twice, and the prototype solves for all components and the residual at
every application: $4\times20\times2\times4=640$ solves). The structured route spends three quarters of its solves on its spectral diagnostic, mostly power iteration
for the leading eigenvalues of every component. Here power iteration is needed only when the bound is not
small. (A bound that is not small can also mean several important directions, which the interval of
Section~\ref{sec:q1} does not address.)

Three of these four points are numerical and would speed up the earlier variance models as well; only the
third is specific to the variance estimates of this note.

\paragraph{What the comparison leaves out.} The counts exclude the point estimate, which needs hundreds of
solves in \fevc{} and, in the prototype of this note, a factorisation (Section~\ref{sec:computation}). The earlier
routes also deliver more for their cost: both include the eigenvalue calculations and the $q=1$ interval of
Section~\ref{sec:q1}, and the split-sample estimator of KSS's published paper, which neither code implements, needs no model
for the variances. (In the $q=1$ case the KSS Matlab package fills the three entries of \eqref{eq:q1cov} with
its smoothed variances. Section~\ref{sec:q1} does the same for two of them with the estimates of this note, and
for the first takes the larger of the leave-out estimate of KSS's paper and the value from the estimates of
this note.) Finally, sixteen draws suffice
because $r_{\rm eff}$ is large. With a dominant direction it is not, the bound of Proposition~\ref{prop:sim}
is weak, and more draws are needed.
